\section{治理与可解释性}

\subsection{Drift Playbook}

\begin{figure}[H]
\centering
\includegraphics[width=0.8\textwidth]{slide-02-08.jpg}
\caption{Slide 02-08 展示了 drift response workflow，包括 alert、human review、rollback。}
\end{figure}

Drift Playbook 包括：Detect（alert rules）、Diagnose（cluster drift prompt）、Decide（governance board）、Document（runbook entry）。每次响应都要附带 snapshot、KL 曲线、human override log。

\begin{table}[H]
\centering
\begin{tabular}{p{0.22\textwidth} p{0.65\textwidth}}
\toprule
阶段 & 核心 artifact \\
\midrule
Detect & Reward drift alert + latency spike dashboard \\
Diagnose & Prompt clustering report + human override video \\
Decide & Governance board minutes + rollback checklist \\
Document & Runbook entry + postmortem draft \\
\bottomrule
\end{tabular}
\caption{Drift Playbook 的 artifact delivery}
\end{table}

\begin{knowledgebox}{自动化报警的条目}
建议用 KL drift、hallucination flag、human override frequency 三线报警；只有三条同时越界时触发 governance review，避免频繁 false positive。
\end{knowledgebox}

\subsection{Automated Evaluation Matrix}

Automation coverage包含 unit-level tests（safety prompt injection, allowed domains）、system-level evaluation（trajectory success rate）、ops-level metrics（human override latency）。这些指标要在 dashboard 上分层展示，并与 benchmark/human eval 做 cross-check。

\begin{importantbox}{Evaluation matrix 的多层次意义}
将 benchmark、safety、ops 分层，再用 color-coded heatmap 表示健康状况，可以让 compliance、ops 和 research 在一次 review 中共享事实感知。
\end{importantbox}

\subsection{人类监督与治理会议}

每个版本都需要指定 alignment owner、ops owner、legal owner。Alignment owner 维护 evidence matrix；ops owner 监控 drift dashboards；legal owner 决定 high-risk prompt 是否需要额外 guardrail。

\begin{warningbox}{不要忽视 governance meeting}
如果 governance meeting 只是走过场，团队就会在指标冲突时选择对齐最小的那条线（通常是 benchmark）。必须在会前准备 evidence scorecard，并在会后发布 action items。
\end{warningbox}

\subsection{本章小结}

治理与可解释性的闭环是一种 accountability practice：高质量的 drift playbook、automated evaluation matrix 与多角色的 governance meeting 让模仿学习的 deployment 具备 traceability 与迅速响应能力。

